% Parte de ejercicios: va después de todos los temas, así que puede remitir a cualquier resultado de los apuntes.
\section{Ejercicios de los boletines}

Aquí están todos los ejercicios de los cinco boletines, con el enunciado completo. Para no repetir nada, la solución se da de una de estas formas, por orden de prioridad:
\begin{itemize}
    \item \etiqueta{teal!70!black}{En los apuntes}: la solución ya está en los temas anteriores (casi siempre porque se hizo en clase), y solo se da el enlace.
    \item \fuenteclase{fecha}: se resolvió en clase y no forma parte de la teoría, así que solo está aquí.
    \item \fuenteexamen{\dots}: la solución publicada por el profesor en un examen de 2025-26.
    \item \fuentepropia: el ejercicio no se resolvió en clase ni en ningún examen.
\end{itemize}
\enexamen{\dots} indica que el ejercicio, o uno casi idéntico, se preguntó en un examen.

Los ejercicios que el profesor \emph{propuso en clase} (``Ejercicio:'' en los contenidos de las clases) no están aquí salvo que también estén en un boletín. Como hacen falta para seguir la teoría, están demostrados en su sitio dentro de los temas y llevan la marca \marcaE.

\inicioparte{e1}{Ejercicios del Tema 1}{Resultados preliminares de topología general}{subsubsection}
\subsection{Tema 1. Resultados preliminares de topología general}
\indicetema{e1}

\begin{ejercicio}{T1.1 (continuidad y cerrados)}
Prueba que, si $f \colon X \to Y$ es una aplicación entre espacios topológicos, entonces $f$ es continua si y solo si $f^{-1}(C)$ es cerrado en $X$, para todo cerrado $C \subseteq Y$.

\enapuntes{\cref{prop:prop_lema_previo} (ejercicio de clase [C1.16]).}
\end{ejercicio}

\begin{ejercicio}{T1.2 (proyección estereográfica)}
Prueba que la proyección estereográfica es un homeomorfismo entre $\mathbb{R}^2$ y $S^2 \setminus \{(0,0,1)\}$, siendo $S^2 \subseteq \mathbb{R}^3$ la esfera de centro el origen y radio $1$.

\enapuntes{\cref{prop:proy_estereografica} (ejercicio de clase [C1.7]; se prueba para $S^n$).}
\end{ejercicio}

\begin{ejercicio}{T1.3 (topología inducida en dos pasos)}
Sea $(X,\tau)$ un espacio topológico y $Z \subseteq Y \subseteq X$. Sean $\tau_Y$ y $\tau_Z$ las topologías inducidas por $\tau$ en $Y$ y $Z$. Prueba que $\tau_Z$ es idéntica a la topología en $Z$ inducida por $\tau_Y$.

\sol \fuentepropia

Denotemos por $(\tau_Y)_Z = \{V \cap Z : V \in \tau_Y\}$ la topología inducida por $\tau_Y$ en $Z$.

$(\tau_Y)_Z \subseteq \tau_Z$: si $W = V \cap Z$ con $V \in \tau_Y$, existe $U \in \tau$ con $V = U \cap Y$. Como $Z \subseteq Y$, se tiene $W = U \cap Y \cap Z = U \cap Z \in \tau_Z$.

$\tau_Z \subseteq (\tau_Y)_Z$: si $W = U \cap Z$ con $U \in \tau$, entonces $W = (U \cap Y) \cap Z$ (de nuevo porque $Z \subseteq Y$), y $U \cap Y \in \tau_Y$. Luego $W \in (\tau_Y)_Z$.
\end{ejercicio}

\begin{ejercicio}{T1.4 (caracterizaciones de la conexión)}
Denotemos por $I_2$ el conjunto $\{0,1\}$ con la topología discreta. Sea $X$ un espacio topológico. Prueba que son equivalentes:
\begin{enumerate}
    \item[a)] $X$ es conexo.
    \item[b)] Los únicos subconjuntos abiertos y cerrados de $X$ son $\emptyset$ y $X$.
    \item[c)] Ninguna aplicación continua $f \colon X \to I_2$ es sobreyectiva.
\end{enumerate}

\sol \fuentepropia

\emph{a) $\Leftrightarrow$ b).} $X$ es disconexo si y solo si existe una desconexión $(A,B)$: $A,B$ abiertos, no vacíos, disjuntos y con $A \cup B = X$. Eso equivale a que $A$ sea abierto y cerrado (su complementario $B$ es abierto) con $A \neq \emptyset$ y $A \neq X$. Es decir, $X$ es disconexo si y solo si existe un abierto y cerrado distinto de $\emptyset$ y de $X$.

\emph{b) $\Rightarrow$ c).} Sea $f \colon X \to I_2$ continua. Como $\{0\}$ es abierto y cerrado en $I_2$ (topología discreta), $A = f^{-1}(0)$ es abierto y cerrado en $X$ (\textbf{T1.1}). Por b), $A = \emptyset$ o $A = X$, es decir, $f$ no toma el valor $0$ o no toma el valor $1$. Luego $f$ no es sobreyectiva.

\emph{c) $\Rightarrow$ b).} Si existiera $A$ abierto y cerrado con $A \neq \emptyset, X$, la aplicación $f(x) = 0$ si $x \in A$ y $f(x) = 1$ si $x \notin A$ sería continua (las antiimágenes de los cuatro subconjuntos de $I_2$ son $\emptyset$, $A$, $X \setminus A$ y $X$, todos abiertos) y sobreyectiva, en contra de c).
\end{ejercicio}

\begin{ejercicio}{T1.5 (unión de conexos con un punto común)}
Aplica el Ejercicio 4 para probar que, si $\{A_\alpha\}_{\alpha \in \Lambda}$ es una familia de subespacios conexos de $X$ con $\bigcap_{\alpha} A_\alpha \neq \emptyset$, entonces $\bigcup_{\alpha} A_\alpha$ es conexo.

\sol \fuentepropia

Sean $A = \bigcup_\alpha A_\alpha$, $p \in \bigcap_\alpha A_\alpha$ y $f \colon A \to I_2$ continua. Para cada $\alpha$, la restricción $f|_{A_\alpha}$ es continua (por \textbf{T1.3}, la topología de $A_\alpha$ como subespacio de $A$ es la misma que como subespacio de $X$). Como $A_\alpha$ es conexo, por T1.4 c) $f|_{A_\alpha}$ no es sobreyectiva, es decir, es constante, y vale $f(p)$ porque $p \in A_\alpha$. Luego $f \equiv f(p)$ en todo $A$, no es sobreyectiva y, por T1.4, $A$ es conexo.
\end{ejercicio}

\begin{ejercicio}{T1.6 (unión de conexos que cortan a uno fijo)}
Sea $\{A_\alpha\}_{\alpha \in \Lambda}$ una familia de subespacios conexos de $X$. Prueba que, si existe $\alpha_0 \in \Lambda$ tal que $A_\alpha \cap A_{\alpha_0} \neq \emptyset$ para todo $\alpha$, entonces $\bigcup_\alpha A_\alpha$ es conexo.

\sol \fuentepropia

Para cada $\alpha$, sea $B_\alpha = A_\alpha \cup A_{\alpha_0}$. Es conexo por \textbf{T1.5}, aplicado a la familia $\{A_\alpha, A_{\alpha_0}\}$, cuya intersección es no vacía. Además $\bigcap_\alpha B_\alpha \supseteq A_{\alpha_0} \neq \emptyset$, luego, por T1.5 otra vez, $\bigcup_\alpha B_\alpha = \bigcup_\alpha A_\alpha$ es conexo.
\end{ejercicio}

\begin{ejercicio}{T1.7 (unión de conexos que se cortan dos a dos)}
Sea $\{A_\alpha\}_{\alpha \in \Lambda}$ una familia de subespacios conexos de $X$. Prueba que, si $A_\alpha \cap A_\beta \neq \emptyset$ para todo $\alpha, \beta \in \Lambda$, entonces $\bigcup_\alpha A_\alpha$ es conexo (aplica el Ejercicio 4).

\sol \fuentepropia

Sea $f \colon \bigcup_\alpha A_\alpha \to I_2$ continua. Como en T1.5, cada $f|_{A_\alpha}$ es constante, igual a cierto $c_\alpha$. Dados $\alpha, \beta$, tomamos $x \in A_\alpha \cap A_\beta$, y así $c_\alpha = f(x) = c_\beta$. Luego $f$ es constante, no es sobreyectiva y, por T1.4, la unión es conexa.
\end{ejercicio}

\begin{ejercicio}{T1.8 (arco-conexo, no localmente arco-conexo)}
Sea $S = \{\frac1n\}_{n \geq 1} \cup \{0\}$ y $X = (S \times [0,1]) \cup ([0,1] \times \{0\}) \subseteq \mathbb{R}^2$. Representa $X$, comprueba que es arco-conexo y que no es localmente arco-conexo (observa la intersección de $X$ con $B((0,\frac12),\frac14)$).

\sol \fuentepropia

$X$ es un ``peine'': la base $[0,1] \times \{0\}$ con dientes verticales $\{s\} \times [0,1]$ en cada $s \in S$.

\emph{Arco-conexo.} Todo punto $(s,y) \in X$ se une con $(s,0)$ por el segmento vertical $\{s\} \times [0,y] \subseteq X$, y $(s,0)$ se une con $(0,0)$ por la base. Por tanto todo punto se une con $(0,0)$ y, concatenando, dos puntos cualesquiera se unen por un arco.

\emph{No localmente arco-conexo en $p = (0,\frac12)$.} Sea $U = B(p,\frac14) \cap X$, entorno de $p$ en $X$. Los puntos de la base distan $\frac12$ de $p$, luego $U \subseteq S \times \left]\frac14,\frac34\right[$. Sea $V \subseteq U$ cualquier entorno de $p$. Como $(\frac1n,\frac12) \to p$, existe $n$ con $q = (\frac1n,\frac12) \in V$. Si $\gamma \colon I \to V$ fuera un arco de $p$ a $q$, su primera coordenada $\gamma_1 \colon I \to S$ sería continua; $\gamma_1(I)$ sería conexo, y los únicos conexos de $S$ son los puntos (entre dos puntos de $S$ siempre hay un real que no está en $S$), luego $\gamma_1$ sería constante. Pero $\gamma_1(0) = 0 \neq \frac1n = \gamma_1(1)$. Así, ningún entorno $V \subseteq U$ de $p$ es arco-conexo.
\end{ejercicio}

\begin{ejercicio}{T1.9 (compacto en Hausdorff es cerrado)}
Prueba que, si $K \subseteq X$ es compacto y $X$ es Hausdorff, entonces $K$ es cerrado.

\sol \fuentepropia

Sea $x \in X \setminus K$. Para cada $y \in K$ existen abiertos disjuntos $U_y \ni y$ y $V_y \ni x$ (dentro de los entornos disjuntos que da la propiedad de Hausdorff). Los $U_y$ recubren $K$; por compacidad, $K \subseteq U_{y_1} \cup \dots \cup U_{y_n}$. Entonces $V = V_{y_1} \cap \dots \cap V_{y_n}$ es un abierto que contiene a $x$ y no corta a $U_{y_1} \cup \dots \cup U_{y_n} \supseteq K$. Luego $X \setminus K$ es entorno de cada uno de sus puntos, es abierto, y $K$ es cerrado.
\end{ejercicio}

\begin{ejercicio}{T1.10 (caracterizaciones de Hausdorff)}
Sea $X$ un espacio topológico. Prueba que son equivalentes:
\begin{enumerate}
    \item[a)] $X$ es Hausdorff.
    \item[b)] Para todo $x \in X$ y todo $y \neq x$, existe $U \in E(x)$ tal que $y \notin \overline{U}$.
    \item[c)] Para todo $x \in X$, $\bigcap_{U \in E(x)} \overline{U} = \{x\}$.
    \item[d)] $\Delta = \{(x,x) : x \in X\}$ es cerrado en $X \times X$.
\end{enumerate}

\sol \fuentepropia

\emph{a) $\Rightarrow$ b).} Sean $U \in E(x)$ y $V \in E(y)$ disjuntos. Entonces $V$ contiene un abierto $W \ni y$ con $W \cap U = \emptyset$, luego $y \notin \overline{U}$.

\emph{b) $\Rightarrow$ a).} Dados $x \neq y$, sea $U \in E(x)$ con $y \notin \overline{U}$. Entonces $V = X \setminus \overline{U}$ es un abierto que contiene a $y$ y $U \cap V = \emptyset$, pues $U \subseteq \overline{U}$.

\emph{b) $\Leftrightarrow$ c).} Siempre $x \in \bigcap_{U \in E(x)} \overline{U}$. La igualdad c) dice que todo $y \neq x$ queda fuera de algún $\overline{U}$, que es exactamente b).

\emph{a) $\Rightarrow$ d).} Sea $(x,y) \notin \Delta$, es decir, $x \neq y$. Tomamos $U \in E(x)$ y $V \in E(y)$ disjuntos. Entonces $U \times V$ es entorno de $(x,y)$ y no corta a $\Delta$: si $(z,z) \in U \times V$, sería $z \in U \cap V$. Luego $(X \times X) \setminus \Delta$ es abierto.

\emph{d) $\Rightarrow$ a).} Si $x \neq y$, $(x,y)$ está en el abierto $(X \times X) \setminus \Delta$. Por la definición de la topología producto, existen $U \in E(x)$ y $V \in E(y)$ con $U \times V \subseteq (X \times X) \setminus \Delta$, y eso significa $U \cap V = \emptyset$.
\end{ejercicio}

\begin{ejercicio}{T1.11 (aplicaciones con llegada Hausdorff)}
Sean $f, g \colon X \to Y$ continuas, con $Y$ Hausdorff. Prueba:
\begin{enumerate}
    \item[a)] $\{x \in X : f(x) = g(x)\}$ es cerrado en $X$.
    \item[b)] Si $D$ es denso en $X$ y $f|_D = g|_D$, entonces $f = g$.
    \item[c)] La gráfica de $f$ es un cerrado de $X \times Y$.
    \item[d)] Si $f$ es inyectiva, entonces $X$ es Hausdorff.
\end{enumerate}

\sol \fuentepropia

\textbf{a)} La aplicación $h = (f,g) \colon X \to Y \times Y$ es continua y $\{f = g\} = h^{-1}(\Delta_Y)$, que es cerrado porque $\Delta_Y$ lo es (\textbf{T1.10} d) y por T1.1.

\textbf{b)} $\{f = g\}$ es un cerrado que contiene a $D$, luego contiene a $\overline{D} = X$.

\textbf{c)} La aplicación $k \colon X \times Y \to Y \times Y$, $k(x,y) = (f(x),y)$, es continua, y la gráfica $\{(x,y) : y = f(x)\}$ es $k^{-1}(\Delta_Y)$, que es cerrado.

\textbf{d)} Si $x \neq x'$, entonces $f(x) \neq f(x')$, y existen abiertos disjuntos $U \ni f(x)$ y $V \ni f(x')$. Así, $f^{-1}(U)$ y $f^{-1}(V)$ son abiertos disjuntos que contienen a $x$ y a $x'$.
\end{ejercicio}

\begin{ejercicio}{T1.12 (biyección continua de compacto en Hausdorff)}
Prueba que si $f \colon X \to Y$ es una biyección continua, $X$ es compacto e $Y$ es Hausdorff, entonces $f$ es un homeomorfismo.

\sol \fuentepropia

Basta ver que $f^{-1}$ es continua, es decir (por \textbf{T1.1}), que $(f^{-1})^{-1}(C) = f(C)$ es cerrado para todo cerrado $C \subseteq X$.
\begin{itemize}
    \item $C$ es compacto: si $\{U_i\}$ es un recubrimiento abierto de $C$, entonces $\{U_i\} \cup \{X \setminus C\}$ recubre $X$, y un subrecubrimiento finito de este da uno finito de $C$.
    \item $f(C)$ es compacto, por ser imagen continua de un compacto [C1.14].
    \item $f(C)$ es cerrado en $Y$, por ser compacto en un Hausdorff (\textbf{T1.9}).
\end{itemize}
\end{ejercicio}

\begin{ejercicio}{T1.13 (lema de Lebesgue)}
Sea $(X,d)$ un espacio métrico compacto y $\mathcal{U} = \{U_\alpha\}_{\alpha \in \Lambda}$ un recubrimiento abierto de $X$. Entonces existe $\delta > 0$ tal que, para todo $x \in X$, existe $\alpha \in \Lambda$ con $B(x,\delta) \subseteq U_\alpha$.

\enapuntes{\cref{thm:lema_lebesgue} [C7.4].}
\end{ejercicio}

\inicioparte{e2}{Ejercicios del Tema 2}{Equivalencia homotópica}{subsubsection}
\subsection{Tema 2. Equivalencia homotópica}
\indicetema{e2}

\begin{ejercicio}{T2.1 (homotopías en un producto)}
Sean $f, g \colon X \to Y \times Z$ continuas, con $f = (f_1,f_2)$ y $g = (g_1,g_2)$. Prueba que $f \simeq g$ si y solo si $f_1 \simeq g_1$ y $f_2 \simeq g_2$.

\sol \fuentepropia

$\Rightarrow$: Sea $F \colon f \simeq g$ y sean $p_Y, p_Z$ las proyecciones (continuas). Entonces $F_1 = p_Y \circ F$ es continua, con $(F_1)_0 = p_Y \circ f = f_1$ y $(F_1)_1 = g_1$, luego $F_1 \colon f_1 \simeq g_1$. Análogamente, $p_Z \circ F \colon f_2 \simeq g_2$.

$\Leftarrow$: Sean $F_1 \colon f_1 \simeq g_1$ y $F_2 \colon f_2 \simeq g_2$. La aplicación $F = (F_1,F_2) \colon X \times I \to Y \times Z$ es continua (sus dos componentes lo son) y cumple $F_0 = (f_1,f_2) = f$ y $F_1 = (g_1,g_2) = g$.
\end{ejercicio}

\begin{ejercicio}{T2.2 (aplicaciones al hemisferio norte)}
Sea $X$ un espacio topológico y sea $Y$ el hemisferio norte abierto de $S^2$ (sin el ecuador). Prueba que cualquier par de aplicaciones continuas $X \to Y$ son homotópicas.

\sol \fuentepropia

Sea $Y = \{(x,y,z) \in S^2 : z > 0\}$ y sean $f,g \colon X \to Y$ continuas. Definimos
\[ F(x,t) = \frac{(1-t)f(x) + t\,g(x)}{\|(1-t)f(x) + t\,g(x)\|}. \]
La tercera coordenada de $(1-t)f(x) + t\,g(x)$ es combinación convexa de dos números positivos, luego es positiva. En particular, el vector no es nulo, $F$ está bien definida y es continua, y $F(x,t)$ tiene norma $1$ y tercera coordenada positiva, así que $F(x,t) \in Y$. Además $F_0 = f$ y $F_1 = g$ (ya tienen norma $1$). Luego $f \simeq g$.

\emph{Otra forma:} $Y \cong \{(x,y) : x^2 + y^2 < 1\}$ mediante $(x,y,z) \mapsto (x,y)$, con inversa $(x,y) \mapsto (x,y,\sqrt{1-x^2-y^2})$. El disco es convexo, luego contráctil; así que $Y$ es contráctil y se aplica \textbf{T2.6} c).
\end{ejercicio}

\begin{ejercicio}{T2.3 (arco-conexo si y solo si $[P,X]$ tiene un elemento)}
Sea $P = \{p\}$ un espacio con un único punto y $X$ un espacio topológico. Prueba que $X$ es arco-conexo $\iff$ el conjunto cociente $[P,X]$ tiene un único elemento.

\sol \fuentepropia

La aplicación $C(P,X) \to X$, $f \mapsto f(p)$, es biyectiva. Veamos que, para $f,g \in C(P,X)$,
\[ f \simeq g \iff \text{existe un arco en } X \text{ de } f(p) \text{ a } g(p). \]
$\Rightarrow$: si $F \colon f \simeq g$, entonces $\alpha(t) = F(p,t)$ es un arco de $f(p)$ a $g(p)$.\\
$\Leftarrow$: si $\alpha$ es un arco de $f(p)$ a $g(p)$, entonces $F(p,t) = \alpha(t)$ es continua (pues $P \times I \cong I$) y $F \colon f \simeq g$.

Por tanto, las clases de $[P,X]$ se corresponden con las arco-componentes de $X$. Así, $[P,X]$ tiene un único elemento si y solo si dos puntos cualesquiera de $X$ se unen por un arco, es decir, si y solo si $X$ es arco-conexo (con $X \neq \emptyset$).
\end{ejercicio}

\begin{ejercicio}{T2.4 (homotópica a la identidad en un discreto)}
Sea $X$ un espacio discreto y $f \colon X \to X$ continua. Prueba que si $f \simeq 1_X$, entonces $f = 1_X$.

\sol \fuentepropia

\emph{Previo: todo arco $\alpha \colon I \to X$ es constante.} $\alpha(I)$ es conexo. En un discreto cada $\{x\}$ es abierto y cerrado, así que un subconjunto con dos puntos o más es disconexo. Luego $\alpha(I)$ es un punto.

Sea $F \colon f \simeq 1_X$. Para cada $x \in X$, $t \mapsto F(x,t)$ es un arco en $X$, luego constante: $f(x) = F(x,0) = F(x,1) = x$. Así $f = 1_X$.
\end{ejercicio}

\begin{ejercicio}{T2.5 (de conexo a discreto)}
Sean $X$ conexo e $Y$ discreto. Prueba que si $f, g \colon X \to Y$ son continuas y homotópicas, entonces $f = g = \text{cte}$.

\sol \fuentepropia

$f(X)$ es un conexo del discreto $Y$, luego es un punto (como en T2.4): $f$ es constante, e igualmente $g$. Si $F \colon f \simeq g$, para cada $x$ el arco $t \mapsto F(x,t)$ es constante (T2.4), así que $f(x) = g(x)$. Luego $f = g$.
\end{ejercicio}

\begin{ejercicio}{T2.6 (caracterización de los contráctiles)}
Prueba que son equivalentes para un espacio topológico $X$:
\begin{enumerate}
    \item[a)] $X$ es contráctil.
    \item[b)] $1_X \simeq c_x$, para todo $x \in X$.
    \item[c)] Para todo par de aplicaciones $f, g \colon T \to X$ se tiene $f \simeq g$.
    \item[d)] $X$ es del mismo tipo de homotopía que un punto.
\end{enumerate}

\enapuntes{\cref{thm:teo_contractil} (ejercicio de clase [C2.12]).} \quad \enexamen{final 29/01/2026, ej.\ 1b (0,5 pts)}

En el examen bastaba con una caracterización: el profesor probó solo a)$\Leftrightarrow$d) y citó c).
\end{ejercicio}

\begin{ejercicio}{T2.7 (contráctil $\Rightarrow$ arco-conexo)}
Prueba que si $X$ es contráctil, entonces $X$ es arco-conexo. (El recíproco no es cierto: $S^1$ no es contráctil.)

\enapuntes{\cref{thm:teo_contractil_arcoconexo} [C3.6]. El recíproco es falso: \cref{obs:obs_recip_contractil} ($S^1$ es arco-conexo y no contráctil, \cref{prop:cor_s1_no_contractil}).} \quad \enexamen{parcial 31/10/2025, ej.\ 1b (incluye el recíproco)}
\end{ejercicio}

\begin{ejercicio}{T2.8 (producto de contráctiles)}
Prueba que $X_1 \times \dots \times X_r$ es contráctil si y solo si $X_i$ es contráctil para todo $i$.

\sol \fuentepropia

Escribimos $X = X_1 \times \dots \times X_r$ (todos no vacíos).

$\Leftarrow$: sean $F_i \colon 1_{X_i} \simeq c_{a_i}$. Entonces
\[ F\big((x_1,\dots,x_r),t\big) = \big(F_1(x_1,t),\dots,F_r(x_r,t)\big) \]
es continua (por componentes), $F_0 = 1_X$ y $F_1 = c_{(a_1,\dots,a_r)}$.

$\Rightarrow$: sea $F \colon 1_X \simeq c_{(a_1,\dots,a_r)}$. Fijado $i$, elegimos $b_j \in X_j$ para $j \neq i$ y sea $\iota_i \colon X_i \to X$ la aplicación que pone $x$ en el lugar $i$ y $b_j$ en los demás. Si $p_i$ es la proyección, $G(x,t) = p_i\big(F(\iota_i(x),t)\big)$ es continua, $G(x,0) = p_i(\iota_i(x)) = x$ y $G(x,1) = a_i$. Luego $1_{X_i} \simeq c_{a_i}$.
\end{ejercicio}

\begin{ejercicio}{T2.9 (el cono es contráctil)}
Sea $C(X) = (X \times [0,1])/\!\sim$, donde se colapsa $X \times \{1\}$ a un punto, con la topología cociente. Prueba que $C(X)$ es contráctil. (Ayuda: $F([(x,s)],t) = [(x,(1-t)s+t)]$.)

\sol \fuentepropia

Sea $q \colon X \times I \to C(X)$ la proyección y $v = q(X \times \{1\})$ el vértice.

\emph{Bien definida.} Si $[(x,s)] = [(y,s')]$ con $(x,s) \neq (y,s')$, entonces $s = s' = 1$, y $(1-t)\cdot 1 + t = 1$, así que ambas imágenes son $v$.

\emph{Continua.} Sea $G \colon X \times I \times I \to X \times I$, $G(x,s,t) = (x,(1-t)s+t)$, que es continua (nótese que $(1-t)s + t \in [0,1]$). Se cumple $F \circ (q \times 1_I) = q \circ G$, que es continua. Como $I$ es compacto Hausdorff, $q \times 1_I$ es una identificación (lema de Whitehead), luego $F$ es continua.

\emph{Extremos.} $F_0([(x,s)]) = [(x,s)]$, luego $F_0 = 1_{C(X)}$; y $F_1([(x,s)]) = [(x,1)] = v$, luego $F_1 = c_v$. Así $C(X)$ es contráctil.
\end{ejercicio}

\begin{ejercicio}{T2.10 (unión de contráctiles)}
Sean $X$ e $Y$ contráctiles con $X \cap Y$ no vacío y arco-conexo. ¿Es $X \cup Y$ contráctil? (Usa que $S^1$ no es contráctil.)

\sol \fuentepropia

\textbf{No.} Dentro de $S^1$, sean
\[ X = \{e^{i\theta} : \theta \in [0,\pi]\}, \qquad Y = \{e^{i\theta} : \theta \in \left]\pi, 2\pi\right]\}. \]
$X \cong [0,\pi]$ e $Y \cong \left]\pi,2\pi\right]$ (mediante $\theta \mapsto e^{i\theta}$, que es inyectiva en cada intervalo; para $Y$ la inversa es continua porque $Y$ no llega a $e^{i\pi}$ y el ángulo se puede tomar en $\left]\pi,2\pi\right]$ de forma continua). Ambos intervalos son convexos, luego contráctiles. Además $X \cap Y = \{(1,0)\}$ es no vacío y arco-conexo. Pero $X \cup Y = S^1$, que no es contráctil [C9.6].

\emph{Trampa:} el ejemplo de las dos semicircunferencias \emph{cerradas} no vale, porque su intersección son dos puntos, que no es arco-conexo. Tampoco sirven dos abiertos: por [C7.6] su unión sería simplemente conexa.
\end{ejercicio}

\begin{ejercicio}{T2.11 (aplicaciones de $S^{n-1}$ homotópicas a constantes)}
Sea $f \colon S^{n-1} \to Y$ continua, $n \geq 1$, y $p \in S^{n-1}$. Prueba que son equivalentes:
\begin{enumerate}
    \item[a)] $f$ es homotópica a una aplicación constante.
    \item[b)] $f$ se extiende a una aplicación continua $D^n \to Y$.
    \item[c)] $f \simeq_{\{p\}} c_{f(p)}$.
\end{enumerate}

\sol \fuentepropia

\emph{a)$\Rightarrow$b).} Por simetría de $\simeq$, sea $H \colon c_{y_0} \simeq f$, es decir, $H(x,0) = y_0$ y $H(x,1) = f(x)$. Definimos $\bar f \colon D^n \to Y$:
\[ \bar f(x) = \begin{cases} y_0 & \text{si } 0 \leq \|x\| \leq \frac12, \\[0.5ex] H\left(\dfrac{x}{\|x\|},\, 2\|x\| - 1\right) & \text{si } \frac12 \leq \|x\| \leq 1. \end{cases} \]
Si $\|x\| = \frac12$, la segunda expresión vale $H(\cdot,0) = y_0$, luego está bien definida. Cada trozo es continuo (en el segundo, $\|x\| \geq \frac12 > 0$) y los dos dominios son cerrados, así que $\bar f$ es continua por el lema de continuidad [C1.17]. Si $\|x\| = 1$, $\bar f(x) = H(x,1) = f(x)$: extiende a $f$.

\emph{b)$\Rightarrow$c).} Sea $\bar f \colon D^n \to Y$ continua con $\bar f|_{S^{n-1}} = f$. Definimos
\[ G(x,t) = \bar f\big((1-t)x + t\,p\big), \qquad x \in S^{n-1},\ t \in I. \]
Está bien definida porque $D^n$ es convexo, y es continua. Además $G_0 = f$, $G_1 = c_{\bar f(p)} = c_{f(p)}$ y $G(p,t) = \bar f(p) = f(p)$ para todo $t$. Luego $G \colon f \simeq_{\{p\}} c_{f(p)}$.

\emph{c)$\Rightarrow$a).} Una homotopía relativa a $\{p\}$ es, en particular, una homotopía.
\end{ejercicio}

\begin{ejercicio}{T2.12 (extensión a $D^n$ con llegada contráctil)}
Como corolario, prueba que si $Y$ es contráctil, entonces toda aplicación continua $f \colon S^{n-1} \to Y$ es la restricción de una aplicación continua $D^n \to Y$.

\sol \fuentepropia

Como $Y$ es contráctil, por \textbf{T2.6} c) toda aplicación $S^{n-1} \to Y$ es homotópica a una constante. Por \textbf{T2.11} a)$\Rightarrow$b), $f$ se extiende a $D^n$.
\end{ejercicio}

\begin{ejercicio}{T2.13 (el peine)}
Sea $S = \{\frac1n : n \geq 1\} \cup \{0\}$ y $X = (S \times [0,1]) \cup ([0,1] \times \{0\})$. Sean $Y = \{0\} \times \left]0,1\right]$, $L = [0,1] \times \{0\}$, $r_x = \{x\} \times \left]0,1\right]$ (así $r_0 = Y$), $\overline{r_x} = \{x\} \times [0,1]$ y $a_n = \frac1n$. Prueba:
\begin{enumerate}
    \item[a)] $\overline{r_{a_n}}$ es cerrado en $X$.
    \item[b)] $X \setminus r_{a_n}$ y $\overline{r_{a_n}}$ son dos cerrados de $X$ cuya unión es $X$.
    \item[c)] $X \setminus Y$ no es cerrado en $X$.
    \item[d)] $X$ es contráctil. En particular, es arco-conexo y $1_X \simeq c_p$ para todo $p \in X$.
    \item[e)] $1_X \simeq_{\{p\}} c_p \iff p \in X \setminus Y$.
\end{enumerate}
(Ayuda: si $1_X \simeq_{\{p\}} c_p$, para todo $U \in E(p)$ existe $V \in E(p)$ tal que todo $q \in V$ se une con $p$ por un arco contenido en $U$.)

\sol \fuentepropia

\textbf{a)} $\overline{r_{a_n}} = \{a_n\} \times [0,1]$ es cerrado en $\mathbb{R}^2$, luego $\overline{r_{a_n}} = \overline{r_{a_n}} \cap X$ es cerrado en $X$ (cerrados relativos [C1.5]).

\textbf{b)} Como $a_n$ es un punto aislado de $S$, existe $\varepsilon > 0$ tal que $]a_n - \varepsilon, a_n + \varepsilon[$ no contiene otro punto de $S$. Entonces $r_{a_n} = X \cap \big(]a_n - \varepsilon, a_n + \varepsilon[ \times ]0,2[\big)$ es abierto en $X$, y su complementario $X \setminus r_{a_n}$ es cerrado. Junto con a): $X = (X \setminus r_{a_n}) \cup \overline{r_{a_n}}$, unión de dos cerrados (se cortan solo en $(a_n,0)$).

\textbf{c)} Los puntos $(\frac1n,\frac12) \in X \setminus Y$ convergen a $(0,\frac12) \in Y$. Así, $X \setminus Y$ no contiene a todos sus puntos de acumulación en $X$ y no es cerrado.

\textbf{d)} Definimos $F \colon X \times I \to X$:
\[ F((x,y),t) = \begin{cases} (x,\,(1-2t)y) & \text{si } t \in [0,\frac12], \\ ((2-2t)x,\,0) & \text{si } t \in [\frac12,1]. \end{cases} \]
Primero se bajan los dientes a la base y después se desliza la base hasta $(0,0)$. No sale de $X$: $(x,(1-2t)y)$ está en el mismo diente que $(x,y)$ (o en la base si $y = 0$), y $((2-2t)x,0) \in L$. En $t = \frac12$ ambas expresiones valen $(x,0)$, y por el lema de continuidad $F$ es continua. $F_0 = 1_X$ y $F_1 = c_{(0,0)}$. Así $X$ es contráctil y, por \textbf{T2.6} y \textbf{T2.7}, $1_X \simeq c_p$ para todo $p$ y $X$ es arco-conexo.

\textbf{e)} \emph{Prueba de la ayuda.} Sea $F \colon 1_X \simeq_{\{p\}} c_p$ y $U \in E(p)$; podemos suponer $U$ abierto. Como $F(p,t) = p \in U$ para todo $t$, por continuidad cada $(p,t)$ tiene un entorno $W_t \times J_t$ con $F(W_t \times J_t) \subseteq U$. Por compacidad de $I$, bastan $J_{t_1},\dots,J_{t_k}$, y $V = W_{t_1} \cap \dots \cap W_{t_k}$ cumple $F(V \times I) \subseteq U$. Para $q \in V$, $t \mapsto F(q,t)$ es un arco en $U$ de $q$ a $p$.

\emph{$\Rightarrow$} (contrarrecíproco). Sea $p = (0,y_0) \in Y$, con $y_0 > 0$, y sea $U = B(p,\frac{y_0}{2}) \cap X$. Como $U$ no corta a la base, $U \subseteq S \times \left]\frac{y_0}{2},\frac{3y_0}{2}\right[$. Todo entorno $V$ de $p$ contiene algún $q = (\frac1n,y_0)$. Un arco en $U$ de $q$ a $p$ tendría la primera coordenada continua con valores en $S$, luego constante (como en T1.8), pero pasa de $\frac1n$ a $0$. Por la ayuda, $1_X \not\simeq_{\{p\}} c_p$.

\emph{$\Leftarrow$.} Caso $p = (x_0,0) \in L$: la homotopía de d), cambiando la segunda fase por $\big((2-2t)x + (2t-1)x_0,\,0\big)$, fija $p$ en todo momento y acaba en $c_p$.

Caso $p = (a_n,y_0)$ en un diente $r_{a_n}$: lo hacemos en dos fases que fijan $p$ y concluimos por la transitividad de $\simeq_{\{p\}}$ [C3.9].
\begin{itemize}
    \item \emph{Fase 1.} En el cerrado $X \setminus r_{a_n}$ aplicamos la contracción del caso anterior hacia $(a_n,0)$, que no sale de $X \setminus r_{a_n}$ y fija $(a_n,0)$. En el cerrado $\overline{r_{a_n}}$ dejamos todo fijo. Coinciden en la intersección $\{(a_n,0)\}$, así que por b) y el lema de continuidad se obtiene $1_X \simeq_{\{p\}} g$, donde $g$ manda $X \setminus r_{a_n}$ a $(a_n,0)$ y es la identidad en $\overline{r_{a_n}}$.
    \item \emph{Fase 2.} $K(q,t) = \big(a_n,\,(1-t)\,g_2(q) + t\,y_0\big)$, donde $g_2$ es la segunda coordenada de $g$. Es continua, se queda en $\overline{r_{a_n}}$, fija $p$ (pues $g(p) = p$), y va de $g$ a $c_p$.
\end{itemize}
\end{ejercicio}

\inicioparte{e3}{Ejercicios del Tema 3}{El grupo fundamental}{subsubsection}
\subsection{Tema 3. El grupo fundamental}
\indicetema{e3}

\begin{ejercicio}{T3.1 (la concatenación respeta $\simeq_{\{0,1\}}$)}
Sean $\alpha, \alpha' \in C_{x_0,x_1}(X)$ y $\beta, \beta' \in C_{x_1,x_2}(X)$ con $\alpha \simeq_{\{0,1\}} \alpha'$ y $\beta \simeq_{\{0,1\}} \beta'$. Prueba que $\alpha * \beta \simeq_{\{0,1\}} \alpha' * \beta'$. Observa que, de hecho, $\alpha * \beta \simeq_{\{0,\frac12,1\}} \alpha' * \beta'$.

\enapuntes{\cref{prop:ej_poincare_arcos}~a) (ejercicio de clase [C5.3]).} \quad \enexamen{final 29/01/2026, ej.\ 4a, versión con lazos (1 pt)}

\emph{Lo de $\{0,\frac12,1\}$} \fuentepropia: la homotopía de esa demostración cumple además $H(\frac12,t) = H^{(1)}(1,t) = x_1$ para todo $t$, así que también fija $s = \frac12$.
\end{ejercicio}

\begin{ejercicio}{T3.2 (inverso de una concatenación)}
Prueba que, si $\alpha$ y $\beta$ son arcos con $\alpha(1) = \beta(0)$, entonces $(\alpha * \beta)^{-1} = \beta^{-1} * \alpha^{-1}$.

\enapuntes{\cref{prop:ej_inverso_concat} (ejercicio de clase [C5.7]).}
\end{ejercicio}

\begin{ejercicio}{T3.3 (Poincaré para arcos)}
Utilizando el mismo argumento que en la prueba del teorema de Poincaré, prueba:
\begin{enumerate}
    \item[a)] Si $\alpha \in C_{x_0,x_1}(X)$, entonces $\alpha * c_{x_1} \simeq_{\{0,1\}} c_{x_0} * \alpha \simeq_{\{0,1\}} \alpha$.
    \item[b)] Si $\alpha \in C_{x_0,x_1}(X)$, entonces $\alpha * \alpha^{-1} \simeq_{\{0,1\}} c_{x_0}$ y $\alpha^{-1} * \alpha \simeq_{\{0,1\}} c_{x_1}$.
    \item[c)] Si $\alpha \in C_{x_0,x_1}$, $\beta \in C_{x_1,x_2}$ y $\gamma \in C_{x_2,x_3}$, entonces $\alpha * (\beta * \gamma) \simeq_{\{0,1\}} (\alpha * \beta) * \gamma$.
\end{enumerate}

\enapuntes{\cref{prop:ej_poincare_arcos} (ejercicio de clase [C5.3]). Correspondencia de apartados: a) del boletín es c) de los apuntes, b) es d) y c) es b).}
\end{ejercicio}

\begin{ejercicio}{T3.4 (arcos homotópicos inducen el mismo cambio de base)}
Sean $w, w' \in C_{x_0,x_1}(X)$. Prueba que, si $w \simeq_{\{0,1\}} w'$, entonces $P_w = P_{w'}$.

\enapuntes{\cref{prop:ej_pw_homotopicos} (ejercicio de clase [C5.12]).}
\end{ejercicio}

\begin{ejercicio}{T3.5 ($\pi_1$ abeliano y cambios de base)}
Sea $X$ arco-conexo. Prueba que $\pi_1(X,x)$ es abeliano si y solo si, para todo par de puntos $x_0, x_1 \in X$ y todo par de arcos $w, w' \in C_{x_0,x_1}(X)$, se cumple $P_w = P_{w'}$.

\sol \fuentepropia

\emph{Cálculo previo.} Sea $\delta = w^{-1} * w' \in C_{x_1}(X)$, un lazo en $x_1$. Para $[\alpha] \in \pi_1(X,x_0)$, metiendo $w * w^{-1} \simeq c_{x_0}$ (\textbf{T3.3} b) y usando \textbf{T3.1} y la asociatividad (T3.3 c):
\[ P_{w'}([\alpha]) = [w'^{-1} * w * w^{-1} * \alpha * w * w^{-1} * w'] = [\delta]^{-1} \cdot P_w([\alpha]) \cdot [\delta], \]
donde se ha usado $[w'^{-1} * w] = [(w^{-1} * w')^{-1}] = [\delta]^{-1}$ (\textbf{T3.2}).

$\Rightarrow$: como $X$ es arco-conexo, todos los $\pi_1(X,x)$ son isomorfos [C5.10], luego $\pi_1(X,x_1)$ es abeliano y $[\delta]^{-1} \cdot P_w([\alpha]) \cdot [\delta] = P_w([\alpha])$. Así $P_{w'} = P_w$.

$\Leftarrow$: sea $x_0 = x_1 = x$, $w = c_x$ y $w' = \gamma$ un lazo cualquiera en $x$. Entonces $P_{c_x} = 1$ (por T3.3 a) y $P_\gamma([\alpha]) = [\gamma]^{-1}[\alpha][\gamma]$. La hipótesis $P_\gamma = P_{c_x}$ da $[\gamma]^{-1}[\alpha][\gamma] = [\alpha]$, es decir, $[\alpha][\gamma] = [\gamma][\alpha]$ para todo $[\alpha], [\gamma]$: $\pi_1(X,x)$ es abeliano.
\end{ejercicio}

\begin{ejercicio}{T3.6 (simplemente conexo y arcos con los mismos extremos)}
Sea $X$ arco-conexo. Prueba que son equivalentes:
\begin{enumerate}
    \item[a)] $\pi_1(X)$ es el grupo trivial.
    \item[b)] Para todo par de puntos $x_0, x_1$ y todo par de arcos $\alpha, \beta$ de $x_0$ a $x_1$, se cumple $\alpha \simeq_{\{0,1\}} \beta$.
\end{enumerate}

\sol \fuentepropia

\emph{b)$\Rightarrow$a).} Dado un lazo $\alpha \in C_{x_0}(X)$, aplicamos b) con $x_1 = x_0$ y $\beta = c_{x_0}$: $[\alpha] = [c_{x_0}]$.

\emph{a)$\Rightarrow$b).} $\alpha * \beta^{-1}$ es un lazo en $x_0$; por a), $\alpha * \beta^{-1} \simeq_{\{0,1\}} c_{x_0}$. Usando \textbf{T3.1} y \textbf{T3.3}:
\[ \alpha \simeq \alpha * c_{x_1} \simeq \alpha * (\beta^{-1} * \beta) \simeq (\alpha * \beta^{-1}) * \beta \simeq c_{x_0} * \beta \simeq \beta, \]
todas relativas a $\{0,1\}$.
\end{ejercicio}

\begin{ejercicio}{T3.7 (grupo fundamental de un producto)}
Sean $X$ e $Y$ espacios topológicos y $(x_0,y_0) \in X \times Y$. Prueba que $\pi_1(X \times Y,(x_0,y_0)) \cong \pi_1(X,x_0) \times \pi_1(Y,y_0)$.

\enapuntes{\cref{prop:ej_pi1_producto} (ejercicio de clase [C9.8]).}

\emph{Consecuencias:} $\pi_1(T^2) = \pi_1(S^1 \times S^1) \cong \mathbb{Z}^2$ y $\pi_1(S^1 \times [0,1]) \cong \mathbb{Z} \times 1 \cong \mathbb{Z}$.
\end{ejercicio}

\begin{ejercicio}{T3.8 (producto por un contráctil)}
Sean $X$ e $Y$ espacios topológicos con $Y$ contráctil. Prueba que $X \times Y \simeq X$.

\sol \fuentepropia

Sea $F \colon 1_Y \simeq c_{y_0}$. Tomamos $f \colon X \times Y \to X$, $f(x,y) = x$, y $g \colon X \to X \times Y$, $g(x) = (x,y_0)$. Entonces $f \circ g = 1_X$ y $g \circ f(x,y) = (x,y_0)$. La aplicación
\[ H((x,y),t) = (x,\, F(y,t)) \]
es continua, con $H_0 = 1_{X \times Y}$ y $H_1 = g \circ f$. Luego $g \circ f \simeq 1$ y $X \times Y \simeq X$.
\end{ejercicio}

\begin{ejercicio}{T3.9 (quitar un punto en una equivalencia homotópica)}
Sea $f \colon X \to Y$ una equivalencia homotópica con $f^{-1}(f(p)) = \{p\}$. ¿Es cierto que la restricción $f| \colon X \setminus \{p\} \to Y \setminus \{f(p)\}$ es una equivalencia homotópica?

\sol \fuentepropia

\textbf{No.} Sea $Y = C = [-1,1] \times \{0\} \cup \{0\} \times [-1,1]$ la cruz y $X = \mathbb{R}$. Sea
\[ f \colon \mathbb{R} \to C, \qquad f(x) = \left(\frac{x}{1+|x|},\, 0\right). \]
$X$ es convexo y $C$ es estrellado respecto del origen, luego ambos son contráctiles. Así, \emph{cualquier} $f$ continua es una equivalencia homotópica: con $g \colon C \to \mathbb{R}$ cualquiera, $g \circ f \simeq 1_X$ y $f \circ g \simeq 1_C$ por \textbf{T2.6} c). Además $f^{-1}(f(0)) = f^{-1}(0,0) = \{0\}$.

Pero $f| \colon \mathbb{R} \setminus \{0\} \to C \setminus \{(0,0)\}$ no es una equivalencia homotópica. Supongamos que $g$ fuera una inversa homotópica y $H \colon f \circ g \simeq 1$ en $C \setminus \{(0,0)\}$. Sea $q = (0,\frac12)$, en el brazo superior. El arco $t \mapsto H(q,t)$ une $f(g(q))$, que está en el eje horizontal (la imagen de $f$), con $q$, que está en el brazo superior. Pero el brazo superior $\{0\} \times \left]0,1\right]$ es una arco-componente de $C \setminus \{(0,0)\}$ que no corta al eje horizontal. Contradicción.

(Con el Tema 5 es inmediato: $\#\pi_0(\mathbb{R} \setminus \{0\}) = 2 \neq 4 = \#\pi_0(C \setminus \{(0,0)\})$. Un ejemplo aún más simple: $f \colon \{p\} \to [0,1]$, $f(p) = 0$, es una equivalencia homotópica, pero $\emptyset \to \left]0,1\right]$ no puede serlo, porque no hay aplicaciones $\left]0,1\right] \to \emptyset$.)
\end{ejercicio}

\begin{ejercicio}{T3.10 (homotópica a constante $\Rightarrow$ $f_*$ nulo)}
Sea $f \colon X \to Y$ continua, $y_0 \in Y$ y $\gamma_{y_0} \colon X \to Y$ la constante $x \mapsto y_0$. Prueba que, si $f \simeq \gamma_{y_0}$, entonces $f_* \colon \pi_1(X,x) \to \pi_1(Y,f(x))$ es el homomorfismo nulo, para todo $x \in X$.

\sol \fuenteexamen{parcial 31/10/2025, ej.\ 3} \quad \enexamen{parcial 31/10/2025, ej.\ 3, literal (1,5 pts)}

Sea $x_0 \in X$ y $F \colon f \simeq \gamma_{y_0}$. Sea $\sigma \colon I \to Y$, $\sigma(t) = F(x_0,t)$, que va de $f(x_0)$ a $y_0$. Por el teorema de las aplicaciones homotópicas [C6.7], el diagrama
\[ \begin{array}{ccc} \pi_1(X,x_0) & \xrightarrow{\ f_{*,x_0}\ } & \pi_1(Y,f(x_0)) \\ & \searrow{\scriptstyle (\gamma_{y_0})_{*,x_0}} & \downarrow{\scriptstyle P_\sigma} \\ & & \pi_1(Y,y_0) \end{array} \]
conmuta, es decir, $P_\sigma \circ f_{*,x_0} = (\gamma_{y_0})_{*,x_0}$.

Para $\alpha \in C_{x_0}(X)$, $(\gamma_{y_0})_{*,x_0}([\alpha]) = [\gamma_{y_0} \circ \alpha] = [c_{y_0}]$, luego $(\gamma_{y_0})_{*,x_0}$ es el morfismo nulo. Por tanto, para todo $[\alpha]$,
\[ P_\sigma\big(f_{*,x_0}([\alpha])\big) = [c_{y_0}], \]
es decir, $f_{*,x_0}([\alpha]) \in \ker(P_\sigma) = \{[c_{f(x_0)}]\}$, porque $P_\sigma$ es isomorfismo. Así $f_{*,x_0}$ es el homomorfismo nulo. (Alternativamente: $f_{*,x_0} = P_\sigma^{-1} \circ (\gamma_{y_0})_{*,x_0}$.)

\textbf{Resultados utilizados} (el examen pide enunciarlos):
\begin{enumerate}
    \item[(a)] \emph{[C6.7]} Sean $f,g \colon X \to Y$ continuas, $F \colon f \simeq g$ y $x_0 \in X$. Sea $\sigma(t) = F(x_0,t)$ y $P_\sigma \colon \pi_1(Y,f(x_0)) \to \pi_1(Y,g(x_0))$ el isomorfismo inducido por $\sigma$. Entonces $g_{*,x_0} = P_\sigma \circ f_{*,x_0}$.
    \item[(b)] \emph{[C5.9]} Si $w$ es un arco de $x_0$ a $x_1$, la aplicación $P_w \colon \pi_1(X,x_0) \to \pi_1(X,x_1)$, $P_w([\alpha]) = [w^{-1} * \alpha * w]$, es un isomorfismo de grupos.
\end{enumerate}

\emph{Error típico:} no basta decir ``$f \circ \alpha \simeq \gamma_{y_0} \circ \alpha$''. Esa homotopía es libre (el punto base recorre $\sigma$) y $\gamma_{y_0} \circ \alpha$ es un lazo en otro punto base.
\end{ejercicio}

\inicioparte{e4}{Ejercicios del Tema 4}{Espacios simplemente conexos, elevaciones y $\pi_1(S^n)$}{subsubsection}
\subsection{Tema 4. Espacios simplemente conexos, elevaciones y $\pi_1(S^n)$}
\indicetema{e4}

\begin{ejercicio}{T4.1 (esferas tangentes)}
Sean $S$ y $S'$ dos esferas del mismo radio en $\mathbb{R}^3$ que se cortan en un único punto (tangentes). Prueba que $S \cup S'$ es simplemente conexo. Obtén la misma conclusión para una familia finita de esferas del mismo radio $\{S_1,\dots,S_r\}$ con $S_i$ y $S_{i+1}$ tangentes para todo $i$ y $S_1 \cap S_r = \emptyset$. (Ayuda: aplica la condición suficiente para ser simplemente conexo.)

\sol \fuenteclase{23/10}

Sea $A = S \cup S'$ y $x_0$ el punto de tangencia. Se aplica [C7.6] con un recubrimiento abierto de $A$ por dos conjuntos. Hay que cuidar la topología: una esfera sola \emph{no} es abierta en $A$. Recordemos que $Y \subseteq A$ es abierto en $A$ cuando para todo $p \in Y$ existe $\varepsilon > 0$ con $B(p,\varepsilon) \cap A \subseteq Y$.

Sean $q \in S \setminus \{x_0\}$ y $q' \in S' \setminus \{x_0\}$. Tomamos
\[ V_1 = S \cup (S' \setminus \{q'\}), \qquad V_2 = (S \setminus \{q\}) \cup S'. \]
Ambos son abiertos en $A$ (su complementario es un punto) y $V_1 \cup V_2 = A$.
\begin{itemize}
    \item $V_1$ es una esfera con un disco abierto pegado por un punto ($S' \setminus \{q'\} \cong \mathbb{R}^2$, contráctil). Encogiendo ese disco hasta $x_0$, $V_1 \simeq S^2$, que es simplemente conexo [C7.7]. Igual con $V_2$.
    \item $V_1 \cap V_2 = (S \setminus \{q\}) \cup (S' \setminus \{q'\})$ son dos copias de $\mathbb{R}^2$ unidas por $x_0$, que es arco-conexo.
    \item $x_0 \in V_1 \cap V_2$.
\end{itemize}
Por [C7.6], $A$ es simplemente conexo.

Para la cadena de $r$ esferas se razona por \textbf{inducción} sobre $r$, añadiendo una esfera tangente en cada paso. La hipótesis $S_1 \cap S_r = \emptyset$ es necesaria: si la cadena se cierra formando un anillo de esferas, el espacio \textbf{no} es simplemente conexo.
\end{ejercicio}

\begin{ejercicio}{T4.2 (grupo fundamental de subespacios)}
Determina el grupo fundamental de los siguientes subespacios (topología inducida):
\begin{enumerate}
    \item[a)] $X = \{(x,y) \in \mathbb{R}^2 : 1 \leq x^2+y^2 \leq 4\}$.
    \item[b)] $X = \mathbb{R}^2 \setminus \{(0,0)\}$.
    \item[c)] $X = \{(x,f(x)) : x \in [a,b]\}$, con $f \colon [a,b] \to \mathbb{R}$ continua.
    \item[d)] $X = \{x^2+y^2 = 1\} \cup \{y^2 = 1\}$.
    \item[e)] $X = S^1 \cup \{(x-2)^2+y^2 \leq 1\}$.
    \item[f)] $X = \{(x,y,z) \in \mathbb{R}^3 : 1 \leq x^2+y^2 \leq 4\}$.
    \item[g)] $X = \{x^2+y^2+z^2 = 9,\ 1 \leq z \leq 2\}$.
    \item[h)] $X = \{(x,\cos x,\sin x) : x \in \mathbb{R}\}$.
    \item[i)] $X = \{(x,\cos y,\sin y) : (x,y) \in \mathbb{R}^2\}$.
    \item[j)] $X = \{x(x-1) = 0\} \cup [0,1] \times \{0\}$.
    \item[k)] $X = \{x^2+y^2 > 1\}$.
    \item[l)] $X = [-1,1] \times [-1,1] \times \{0\} \cup \{0\} \times \{0\} \times [-1,1]$.
    \item[m)] $X = \{x \geq 0\} \setminus \{(1,0)\}$.
    \item[n)] $X = \{y \geq 0\} \setminus \{(1,0)\}$.
\end{enumerate}

\textbf{Método.} En cada apartado se busca un $Y$ del mismo tipo de homotopía que $X$ con $\pi_1(Y)$ conocido, y se aplica la invariancia homotópica [C6.8]. Aunque el enunciado no lo pide, conviene dar la equivalencia homotópica explícita, escrita en cartesianas y comprobando que la homotopía \emph{no sale de $X$}.

\medskip
\fuenteclase{23/10, apartados a)--d)}

\textbf{a)} $X \simeq S^1$, luego $\pi_1(X) \cong \mathbb{Z}$. Equivalencia: $f \colon S^1 \to X$ la inclusión y $g \colon X \to S^1$, $g(v) = \frac{v}{\|v\|}$. Se tiene $g \circ f = 1_{S^1}$, y
\[ F(v,t) = (1-t)\,v + t\,\frac{v}{\|v\|} = \big((1-t)\|v\| + t\big)\frac{v}{\|v\|} \]
es continua, con $F_0 = 1_X$ y $F_1 = f \circ g$. No sale de $X$, porque el radio $(1-t)\|v\| + t$ está entre $1$ y $\|v\| \leq 2$. En polares es ``radio $r \mapsto (1-t)r + t$ con el ángulo fijo''; escrito así no hace falta que el ángulo sea continuo, porque $(\cos\theta,\sin\theta) = v/\|v\|$.

\textbf{b)} $\pi_1(X) \cong \pi_1(S^1) \cong \mathbb{Z}$. (Misma $F$ que en a); el coeficiente $(1-t) + \frac{t}{\|v\|}$ es positivo, así que $F(v,t) \neq 0$.)

\textbf{c)} $X$ es \textbf{contráctil} (homeomorfo a $[a,b]$ mediante $x \mapsto (x,f(x))$, con inversa la primera proyección), luego simplemente conexo: $\pi_1(X) = 1$.

\textbf{d)} $X$ es la circunferencia con las dos rectas horizontales $y = \pm 1$, tangentes a ella. $\pi_1(X) \cong \pi_1(S^1) \cong \mathbb{Z}$.

\emph{Equivalencia explícita} \fuentepropia: $r \colon X \to S^1$ es la identidad en $S^1$ y lleva $(x,1) \mapsto (0,1)$ y $(x,-1) \mapsto (0,-1)$. Es continua por el lema de continuidad ($S^1$, $\{y = 1\}$ e $\{y = -1\}$ son tres cerrados que recubren $X$ y las definiciones coinciden en los puntos de tangencia). Con $i$ la inclusión, $r \circ i = 1_{S^1}$, y
\[ H(v,t) = \begin{cases} v & \text{si } v \in S^1, \\ ((1-t)x,\, \pm 1) & \text{si } v = (x,\pm 1) \end{cases} \]
cumple $H_0 = 1_X$ y $H_1 = i \circ r$: encoge cada recta hasta su punto de tangencia.

\medskip
\fuentepropia\ \ (apartados e)--n))

\textbf{e)} Circunferencia más un disco cerrado tangente en $(1,0)$. Sea $D$ el disco. $r \colon X \to S^1$ es la identidad en $S^1$ y vale $(1,0)$ en $D$. $H(v,t) = v$ en $S^1$ y $H(v,t) = (1-t)v + t(1,0)$ en $D$ (convexo, contiene a $(1,0)$); en $(1,0)$ coinciden. Así $X \simeq S^1$ y $\pi_1(X) \cong \mathbb{Z}$.

\textbf{f)} Corona $\times\,\mathbb{R}$. $g(x,y,z) = \frac{(x,y)}{\|(x,y)\|}$ y $f(x,y) = (x,y,0)$, con $g \circ f = 1_{S^1}$. Con $v = (x,y)$:
\[ H((v,z),t) = \left( \big((1-t)\|v\| + t\big)\frac{v}{\|v\|},\ (1-t)z \right), \]
que no sale de $X$ (radio entre $1$ y $2$). $H_0 = 1$ y $H_1 = f \circ g$. Así $\pi_1(X) \cong \mathbb{Z}$.

\textbf{g)} Banda de la esfera de radio $3$ entre $z = 1$ y $z = 2$. En $X$, $x^2+y^2 = 9 - z^2 \in [5,8]$, que no se anula. Sean $g(x,y,z) = \frac{(x,y)}{\|(x,y)\|}$ y $f(x,y) = (\sqrt8\,x,\sqrt8\,y,1)$, con $g \circ f = 1_{S^1}$. Con $z_t = (1-t)z + t \in [1,2]$:
\[ H((x,y,z),t) = \left( \sqrt{9 - z_t^2}\ \frac{(x,y)}{\|(x,y)\|},\ z_t \right) \in X. \]
$H_0 = 1_X$ (porque $\|(x,y)\| = \sqrt{9-z^2}$) y $H_1 = f \circ g$. Así $\pi_1(X) \cong \mathbb{Z}$.

\textbf{h)} Hélice. $\mathbb{R} \to X$, $x \mapsto (x,\cos x,\sin x)$, es continua y biyectiva, con inversa continua la primera proyección. Luego $X \cong \mathbb{R}$, contráctil: $\pi_1(X) = 1$.

\textbf{i)} Cilindro $\mathbb{R} \times S^1$. $g(x,u,v) = (u,v)$ y $f(u,v) = (0,u,v)$, con $g \circ f = 1$. $H((x,u,v),t) = ((1-t)x,u,v)$ va de $1_X$ a $f \circ g$. Así $\pi_1(X) \cong \mathbb{Z}$ (o, por \textbf{T3.7}, $\pi_1(\mathbb{R}) \times \pi_1(S^1) \cong \mathbb{Z}$).

\textbf{j)} Forma de ``H'': las rectas $x = 0$ y $x = 1$ unidas por el segmento $[0,1] \times \{0\}$. Contráctil:
\[ F((x,y),t) = \begin{cases} (x,(1-2t)y) & \text{si } t \in [0,\frac12] \quad \text{(bajar las rectas al segmento)}, \\ ((2-2t)x,0) & \text{si } t \in [\frac12,1] \quad \text{(encoger el segmento)}. \end{cases} \]
No sale de $X$ (las rectas son verticales y el segmento está en $y = 0$), y en $t = \frac12$ ambas valen $(x,0)$. $F_0 = 1$ y $F_1 = c_{(0,0)}$. Así $\pi_1(X) = 1$.

\textbf{k)} Exterior abierto del disco unidad. $S^1 \not\subseteq X$, así que se usa la circunferencia $C_2$ de radio $2$: $g(v) = \frac{2v}{\|v\|}$ y
\[ H(v,t) = \big((1-t)\|v\| + 2t\big)\frac{v}{\|v\|}, \]
cuyo radio está entre $\|v\|$ y $2$, luego es $> 1$. Así $X \simeq C_2 \cong S^1$ y $\pi_1(X) \cong \mathbb{Z}$.

\textbf{l)} Cuadrado más un segmento vertical por el centro. Es \textbf{estrellado} respecto del origen: el segmento de $0$ a un punto del cuadrado está en el cuadrado (convexo), y el de $0$ a un punto del segmento vertical está en ese segmento. Luego $\pi_1(X) = 1$ [C6.5].

\textbf{m)} Semiplano $x \geq 0$ sin $c = (1,0)$. La circunferencia $C$ de centro $c$ y radio $\frac12$ está en $X$ (pues $x \geq \frac12$). Sean $g(v) = c + \frac12 \frac{v-c}{\|v-c\|}$ y, con $\rho = \|v - c\|$ y $u = \frac{v-c}{\rho}$,
\[ H(v,t) = c + \big((1-t)\rho + \tfrac t2\big)\,u. \]
El radio es positivo, así que $H(v,t) \neq c$. La primera coordenada es $1 + \rho_t u_1$, con $\rho_t$ entre $\rho$ y $\frac12$. Si $u_1 \geq 0$ es $\geq 1$; si $u_1 < 0$ es $\geq \min(1 + \rho u_1,\ 1 + \frac{u_1}{2}) \geq \min(x_v, \frac12) \geq 0$, donde $x_v = 1 + \rho u_1 \geq 0$ es la abscisa de $v$. Luego $H$ no sale de $X$, $X \simeq C \cong S^1$ y $\pi_1(X) \cong \mathbb{Z}$.

\textbf{n)} Semiplano $y \geq 0$ sin $(1,0)$. Es \textbf{estrellado} respecto de $(1,1)$: el punto $(1-t)(1,1) + t(x,y)$ tiene ordenada $(1-t) + ty \geq 0$, y vale $(1,0)$ solo si $(1-t) + ty = 0$, lo que obliga a $t = 1$ e $y = 0$, es decir, $(x,y) = (1,0) \notin X$. Luego $\pi_1(X) = 1$.
\end{ejercicio}

\begin{ejercicio}{T4.3 (verdadero o falso)}
\begin{enumerate}
    \item[a)] Existe un recubrimiento abierto $\{U,V\}$ del toro $S^1 \times S^1$ con $U \cap V$ arco-conexo y $U, V$ simplemente conexos.
    \item[b)] Si $\{U,V\}$ es un recubrimiento abierto del cilindro $S^1 \times [0,1]$ con $U$ y $V$ simplemente conexos, entonces $U \cap V$ tiene al menos dos componentes arco-conexas.
    \item[c)] Todo subespacio convexo $X \subseteq \mathbb{R}^n$ es simplemente conexo.
\end{enumerate}

\sol \fuentepropia\ \ (b) coincide con \textbf{T4.4} a), cuya solución es la del profesor)

\textbf{a) Falso.} Si existiera, $U \cap V$ sería no vacío (es arco-conexo, y además el toro es conexo) y arco-conexo, y $U$, $V$ serían simplemente conexos. Por [C7.6], el toro sería simplemente conexo. Pero $\pi_1(S^1 \times S^1) \cong \mathbb{Z} \times \mathbb{Z} \neq 1$ (\textbf{T3.7}).

\textbf{b) Verdadero.} Es \textbf{T4.4} a).

\textbf{c) Verdadero} (para $X \neq \emptyset$). Convexo $\Rightarrow$ estrellado, luego arco-conexo y con $\pi_1(X) = 1$ [C6.5].
\end{ejercicio}

\begin{ejercicio}{T4.4 (verdadero o falso)}
\begin{enumerate}
    \item[a)] Sea $X = S^1 \times [0,1]$ y $\{U,V\}$ un recubrimiento abierto con $U$ y $V$ simplemente conexos. Entonces $U \cap V$ tiene al menos dos componentes arco-conexas.
    \item[b)] Sea $\{V_i\}_{i \in \Lambda}$ un recubrimiento abierto de $X$ con $V_i$ simplemente conexo, $V_i \cap V_j$ arco-conexo para $i \neq j$, y la intersección de todos vacía. Entonces $X$ es simplemente conexo.
    \item[c)] $S^2 \setminus \{(1,0,0),(0,0,1)\}$ es simplemente conexo.
\end{enumerate}

\sol \fuenteexamen{final 29/01/2026, ej.\ 7} \quad \enexamen{final 29/01/2026, ej.\ 7, literal (1,5 pts)}

\textbf{a) Verdadera.} La prueba se basa en [C7.6]: \emph{sea $\{V_i\}$ un recubrimiento abierto de $X$ con $\bigcap_i V_i \neq \emptyset$, $V_i \cap V_j$ arco-conexo para $i \neq j$ y cada $V_i$ simplemente conexo; entonces $X$ es simplemente conexo.}

Supongamos que $U \cap V$ tiene una sola componente arco-conexa. Obviamente $U \cap V \neq \emptyset$, ya que $X$ es conexo. Por el teorema, $X$ sería simplemente conexo. Pero
\[ \pi_1(S^1 \times [0,1]) \cong \pi_1(S^1) \times \pi_1([0,1]) \cong \mathbb{Z}. \]
Contradicción; luego $U \cap V$ tiene al menos dos componentes.

\textbf{b) Falsa.} Que falle la hipótesis $\bigcap V_i \neq \emptyset$ no basta para justificarlo: hace falta un contraejemplo. Sea $X = S^1$, $\epsilon \in \left]0,\frac{1}{100}\right[$ y
\[ V_1 = \{e^{2\pi i\theta} : \theta \in \left]-\epsilon,\tfrac13+\epsilon\right[\},\quad V_2 = \{e^{2\pi i\theta} : \theta \in \left]\tfrac13-\epsilon,\tfrac23+\epsilon\right[\},\quad V_3 = \{e^{2\pi i\theta} : \theta \in \left]\tfrac23-\epsilon,1+\epsilon\right[\}. \]
Cada $V_i$ es abierto en $S^1$ y homeomorfo a un intervalo abierto, luego simplemente conexo. Las intersecciones $V_1 \cap V_2$, $V_1 \cap V_3$ y $V_2 \cap V_3$ son arcos (arco-conexas) y $V_1 \cap V_2 \cap V_3 = \emptyset$. Sin embargo, $\pi_1(S^1) \cong \mathbb{Z}$.

\emph{Observación del profesor:} con las mismas hipótesis también hay ejemplos simplemente conexos (por ejemplo $S^2$ con un recubrimiento análogo), así que sin la hipótesis no se puede concluir nada.

\textbf{c) Falsa.} Sea $X = S^2 \setminus \{(1,0,0),(0,0,1)\}$. Veamos que $X \simeq S^1$, de modo que $\pi_1(X) \cong \mathbb{Z}$. Sea $Y = S^2 \setminus \{(0,0,-1),(0,0,1)\}$; como $X$ e $Y$ son homeomorfos, basta probar $Y \simeq S^1$. Sean
\[ f \colon Y \to S^1,\ f(x,y,z) = \frac{(x,y)}{\|(x,y)\|}, \qquad g \colon S^1 \to Y,\ g(x,y) = (x,y,0). \]
Se tiene $f \circ g = 1_{S^1}$. La aplicación
\[ H((x,y,z),t) = \frac{(x,y,tz)}{\|(x,y,tz)\|} \]
es continua, con $H_0 = g \circ f$ y $H_1 = 1_Y$. Luego $Y \simeq S^1$.

\emph{Complemento} \fuentepropia: la proyección estereográfica desde $(0,0,1)$ da directamente $X \cong \mathbb{R}^2 \setminus \{\text{un punto}\} \simeq S^1$, sin pasar por $Y$.
\end{ejercicio}

\begin{ejercicio}{T4.5 (esfera sin los polos y el ecuador)}
Sea $X = S^2 \setminus \{N,S\}$, con $N = (0,0,1)$ y $S = (0,0,-1)$, y sea $E$ el ecuador. Decide si es cierto o falso:
\begin{enumerate}
    \item[a)] Si $\{U,V\}$ es un recubrimiento abierto de $X$ con $U$ y $V$ simplemente conexos, entonces $U \cap V$ tiene al menos dos componentes conexas.
    \item[b)] Si $f \colon S^2 \to E$ es continua, entonces $f|_E$ no puede ser igual a $1_E$.
    \item[c)] Si $f \colon X \to E$ es continua, entonces $f|_E$ no puede ser igual a $1_E$.
\end{enumerate}

\sol \fuenteexamen{parcial 31/10/2025, ej.\ 5} \quad \enexamen{parcial 31/10/2025, ej.\ 5, literal (2 pts)}

\textbf{a) Verdadera.} $U$ y $V$ son simplemente conexos, luego arco-conexos y no vacíos. Además $U \cap V \neq \emptyset$: si no, $\{U,V\}$ sería una desconexión de $X$, que es conexo. Si $U \cap V$ fuera arco-conexo, [C7.6] daría que $X$ es simplemente conexo. Pero $X \simeq S^1$ (se prueba abajo), que no es simplemente conexo. Luego $U \cap V$ no es arco-conexo: tiene al menos dos componentes arco-conexas. Como $U \cap V$ es abierto y $X$ es localmente arco-conexo, sus componentes conexas coinciden con sus componentes arco-conexas, y $U \cap V$ tiene al menos dos componentes conexas.

\emph{$X \simeq S^1$:} $f(x,y,z) = \frac{(x,y)}{\|(x,y)\|}$ y $g(x,y) = (x,y,0)$ cumplen $f \circ g = 1_{S^1}$, y $H((x,y,z),t) = \frac{(x,y,tz)}{\|(x,y,tz)\|}$ es continua, con $H_0 = g \circ f$ y $H_1 = 1_X$.

\textbf{b) Verdadera.} Supongamos que $f \colon S^2 \to E$ es continua con $f|_E = 1_E$, y sea $j \colon E \to S^2$ la inclusión, de modo que $f \circ j = 1_E$. Para $p_0 \in E$, por funtorialidad,
\[ f_{*,p_0} \circ j_{*,p_0} = (f \circ j)_{*,p_0} = (1_E)_{*,p_0} = 1_{\pi_1(E,p_0)}. \]
Como $\pi_1(S^2,p_0)$ es trivial, la imagen de $f_{*,p_0}$ es el subgrupo trivial; pero $1_{\pi_1(E,p_0)}$ es sobreyectiva sobre $\pi_1(E,p_0) \cong \mathbb{Z}$ ($E \cong S^1$). Contradicción.

\textbf{c) Falsa.} $f(x,y,z) = \frac{(x,y,0)}{\|(x,y,0)\|}$ es continua en $X$ (allí $(x,y) \neq 0$) y su restricción a $E$ es $1_E$.
\end{ejercicio}

\inicioparte{e5}{Ejercicios del Tema 5}{Aplicaciones de la teoría de homotopía}{subsubsection}
\subsection{Tema 5. Aplicaciones de la teoría de homotopía}
\indicetema{e5}

Todos los ejercicios de este boletín se resolvieron en clase: los ejercicios 1, 3 y 4 el 21/10 y los ejercicios 2, 5, 6 y 7 el 23/10. Notación: $p \colon \mathbb{R} \to S^1$, $p(t) = e^{2\pi i t}$, y $\deg_{x_0}(\alpha) = \tilde\alpha(1) - \tilde\alpha(0)$ para \emph{cualquier} elevación $\tilde\alpha$ del lazo $\alpha$ [C10.7]. No depende de la elevación porque dos elevaciones del mismo arco difieren en un entero constante [C10.6].

\begin{ejercicio}{T5.1 (elevación de $w^{-1} * \alpha * w$)}
Sea $x_0 \in S^1$, sea $w \colon I \to S^1$ un arco de $x_0$ a $(1,0)$ y sea $\alpha \in C_{x_0}(S^1)$. Escribe una elevación de $w^{-1} * \alpha * w$ como concatenación de elevaciones de $w^{-1}$, $\alpha$ y $w$. ¿Podemos escribir una elevación de $w$ como el recorrido inverso de una elevación de $w^{-1}$?

\enapuntes{demostración del \cref{thm:teo_cambio_base_grado} (resuelto en clase el 21/10 y el 23/10).} La elevación es $\widetilde{w^{-1}} * \tilde\alpha * \tilde w$, con cada trozo empezando donde acaba el anterior (primer párrafo de la prueba). La respuesta a la pregunta es \textbf{sí}: $(\widetilde{w^{-1}})^{-1}$ eleva $w$, porque $p\big(\widetilde{w^{-1}}(1-t)\big) = w^{-1}(1-t) = w(t)$ (segundo párrafo de la prueba).
\end{ejercicio}

\begin{ejercicio}{T5.2 (el grado no depende del punto base)}
Sea $x_0 \in S^1$ y $w$ un arco de $x_0$ a $(1,0)$. Prueba que $\deg(\alpha) = \deg(w^{-1} * \alpha * w)$ para todo $\alpha \in C_{x_0}(S^1)$. Deduce que, si $\alpha, \beta \in C_{x_0}(S^1)$ y $\alpha \simeq_{\{0,1\}} \beta$, entonces $\deg(\alpha) = \deg(\beta)$.

\enapuntes{\cref{thm:teo_cambio_base_grado} (la igualdad de grados) y \cref{obs:obs_grado_clase} (la consecuencia). Ejercicios de clase [C10.8] y [C10.9], resueltos el 23/10.}
\end{ejercicio}

\begin{ejercicio}{T5.3 (¿Borsuk en el toro?)}
Sea $T^2$ el toro de revolución obtenido al girar alrededor del eje $z$ la circunferencia del plano $x = 0$ de centro $(0,2,0)$ y radio $1$. Dada $f \colon T^2 \to \mathbb{R}^2$ continua, ¿existe siempre un par de puntos opuestos de $T^2$ con la misma imagen?

\sol \fuenteclase{21/10}

\textbf{No.} El toro está centrado en el origen, así que si $p \in T^2$ también $-p \in T^2$. Sea $g \colon \mathbb{R}^3 \to \mathbb{R}^2$, $g(x,y,z) = (x,y)$, y $f = g|_{T^2}$. Si existiera $p = (p_1,p_2,p_3) \in T^2$ con $f(p) = f(-p)$, sería $(p_1,p_2) = (-p_1,-p_2)$, es decir, $p_1 = p_2 = 0$. Pero el toro no tiene puntos con las dos primeras coordenadas nulas (en él, $1 \leq \sqrt{x^2+y^2} \leq 3$). Luego esta $f$ no cumple la conclusión.

No contradice el teorema de Borsuk, que es para $S^2$, y $T^2 \not\cong S^2$ [C9.7].
\end{ejercicio}

\begin{ejercicio}{T5.4 (Borsuk en dimensión 1)}
Sea $f \colon S^1 \to \mathbb{R}$ continua. Prueba que existe $x \in S^1$ con $f(x) = f(-x)$.

\sol \fuenteclase{21/10}

Supongamos que $f(x) \neq f(-x)$ para todo $x \in S^1$. Definimos la aplicación continua
\[ g \colon S^1 \to \mathbb{R}, \qquad g(x) = \frac{f(x) - f(-x)}{|f(x) - f(-x)|}, \]
cuya imagen está contenida en $\{-1,1\}$. Como $S^1$ es conexo, $g(S^1)$ es conexo y contenido en $\{-1,1\}$, luego $g$ es constante: $g \equiv k \in \{-1,1\}$. Pero
\[ g(-x) = \frac{f(-x) - f(x)}{|f(-x) - f(x)|} = -g(x), \]
es decir, $k = -k$, luego $k = 0$, lo cual es absurdo.

\emph{Variante equivalente} \fuentepropia: $h(x) = f(x) - f(-x)$ cumple $h(-x) = -h(x)$; si $h(x_0) \neq 0$, toma un valor positivo y otro negativo, y como $h(S^1)$ es un intervalo, contiene el $0$.
\end{ejercicio}

\begin{ejercicio}{T5.5 ($\pi_0$ es funtorial)}
Para $f \colon X \to Y$ continua, sea $\pi_0(f) \colon \pi_0(X) \to \pi_0(Y)$, $\pi_0(f)(C) = D_C$, la única arco-componente de $Y$ con $f(C) \subseteq D_C$. Prueba:
\begin{enumerate}
    \item[a)] $\pi_0(1_X) = 1_{\pi_0(X)}$.
    \item[b)] $\pi_0(g \circ f) = \pi_0(g) \circ \pi_0(f)$, para $f \colon X \to Y$ y $g \colon Y \to Z$ continuas.
\end{enumerate}

\enapuntes{puntos 1 y 2 de la demostración de la \cref{prop:prop_pi0_invariante} (resuelto en clase el 23/10).}
\end{ejercicio}

\begin{ejercicio}{T5.6 ($\#\pi_0$ es invariante homotópico)}
Prueba que si $X$ e $Y$ son del mismo tipo de homotopía, entonces $\#\pi_0(X) = \#\pi_0(Y)$.

\enapuntes{\cref{prop:prop_pi0_invariante} (resuelto en clase el 23/10).}
\end{ejercicio}

\begin{ejercicio}{T5.7 (seno del topólogo)}
Sea $X = \{(x,\sin\frac1x) : x \in \left]0,1\right]\} \cup (\{0\} \times [-1,1])$. ¿Es $X$ contráctil?

\sol \fuenteclase{23/10}

\textbf{No.} $X$ es la adherencia $\overline{A}$ de la gráfica $A$, el ejemplo clásico de espacio conexo que no es arco-conexo. Tiene $2$ arco-componentes: la gráfica y el segmento $\{0\} \times [-1,1]$. Por tanto $\#\pi_0(X) = 2 \neq 1 = \#\pi_0(\{p\})$. Si $X$ fuese contráctil, tendría el mismo tipo de homotopía que un punto y, por \textbf{T5.6}, $\#\pi_0(X) = 1$. Luego $X$ no es contráctil.

(En clase, los hechos sobre las arco-componentes de $X$ se enunciaron sin demostración.)

\medskip
\emph{Complemento: por qué $X$ no es arco-conexo} \fuentepropia. Basta esto, porque contráctil $\Rightarrow$ arco-conexo (\textbf{T2.7}). Supongamos que existe un arco $\gamma = (\gamma_1,\gamma_2) \colon I \to X$ de $(0,0)$ a $(1,\sin 1)$. El conjunto $\{t : \gamma_1(t) = 0\}$ es cerrado y no vacío; sea $t_0$ su máximo, que cumple $t_0 < 1$. Para $t > t_0$ se tiene $\gamma_1(t) > 0$, luego $\gamma(t)$ está en la gráfica y $\gamma_2(t) = \sin\frac{1}{\gamma_1(t)}$. Como $\gamma_1$ es continua con $\gamma_1(t_0) = 0$, para cada $\delta > 0$ el conjunto $\gamma_1(]t_0,t_0+\delta[)$ es un intervalo de $]0,1]$ con ínfimo $0$, luego contiene algún $]0,\varepsilon[$, y por tanto contiene puntos $\frac{1}{\pi/2 + 2\pi n}$ y $\frac{1}{3\pi/2 + 2\pi n}$ para $n$ grande. Así, $\gamma_2$ toma los valores $1$ y $-1$ en instantes arbitrariamente cercanos a $t_0$, y no puede ser continua en $t_0$. Contradicción.
\end{ejercicio}

